\chapter{考察}
\section{Metric Outdoor事前学習からの適応}
BからCのMAEは13.2730 mから1.1324 m、RMSEは14.5340 mから1.7490 m、AbsRelは0.65896から0.05208へ低下した。BとCはMetric Outdoor由来で出力cap80 mを維持し、同じDay047 validationと65 m以下のGT maskで比較されている。このvalidationにおいて、1 epochの農作物データ適応後に大きな絶対深度改善が観測されたと結論できる。Dは同じ未学習Metric Outdoor初期値から独立に始め、epoch 5のRMSEは1.0094 mであった。ただしCとDはscheduler総期間とrun履歴が異なるため、この差を純粋なepoch数の効果とは断定しない。
\section{絶対深度と局所形状の非同値性}
SiLogは対数深度残差を通して深度場を最適化するが、14 px差分誤差を直接最小化する目的ではない。絶対値や低周波構造が合っても、葉縁等の高周波形状は平滑化され得る。Dの全14 px MAEはepoch 1--5で0.06537から0.06148 mへ改善した一方、GT差20 cm以上の振幅比は0.3046--0.3428で変動し、epoch 5では0.3304であった。absolute depth改善と局所振幅改善は同じ速度・方向ではない。
epoch 2はMAE/RMSE/AbsRelがepoch 1より悪化したが、Pearson、局所14 MAE、振幅比、符号一致は改善した。重み更新、validation分布、最適化状態のどれが原因か保存ログからは特定できない。短いrunにおける非単調推移として扱う。
\section{ゼロ差分基準の含意}
ゼロ差分基準0.06421 mは、GT差分が平坦なペアの多い14 px集計で、常に変化なしとする予測が小さいlocal MAEを得る場合を示す。Bの振幅比は0.3083で、全局所MAE0.06561 mもゼロ基準に近い。これはBの形状が忠実という証拠ではなく、評価平均が振幅縮小に利点を与えることと整合する。GT差0--5 cmが標本の72.96\%を占める一方、20 cm以上は7.48\%である。GT階級、振幅、符号、相関、複数間隔を併記しなければ大きい構造の復元状態を見落とす。
CはGT差20 cm以上でBよりMAEが小さく、符号一致も高いが、振幅比0.327は過小推定を示す。D epoch 5も0.3304であり、課題は残る。第2回と第3回の標本抽出件数・抽出署名の完全な一致は根拠資料から照合できず、CとDの差を対応のある効果量として解釈しない。
\section{Relative由来Aの画像端偽勾配}
第2回では初期head出力は約32.50 m近傍で平坦、学習後に端で強い空間差が見られ、Corn Nadir例ではGTがほぼ一定の路面に対して予測断面が約29.5 mから20 mへ傾く観察があった。これは画像端偽勾配という現象と整合する。しかしzero padding、学習cropと検証時landscape形状の差、新Metric headの学習不足、周辺文脈不足を個別に操作していないので、原因を確定できない。
Aの全画面局所14 MAEは0.10545 mでB/Cより大きいが、端から512 px以上ではA 0.05814 m、B 0.05945 m、C 0.06220 mとなり順位が変わった。境界を除いた値は診断指標であり、主評価へ置き換えていない。Aを全面的に形状不良とすることも内部値だけで優位とすることも適切でない。
\section{518 px入力と写真測量GT}
短辺3956から518 pxへresizeすると元14 pxは入力約1.83 pxとなり、ViT-S/14 patchは元画像換算で約107 pxに相当する。DPTのmulti-scale fusionがpatchより小さい情報を完全に失うという意味ではないが、縮小・patch化・復元補間が葉・茎境界を弱める可能性がある。第2回のGT縮小・復元対照の14 px MAEは0.03117 m、相関0.77472、回帰傾き0.66885であった。これは解像度変換で局所形状が変化し得る診断であり、モデル誤差の説明割合や厳密な誤差下限ではない。
GTは写真測量再構成深度であり、地上実測値と同一ではない。遮蔽、対応点密度、再構成平滑化、位置合わせの不確かさを含み得る。孤立欠損近傍除外で主結果がほぼ変わらなかったことは、系統的なMVS偏りや植生再構成誤差を否定しない。
\section{Oblique画像、作物差、稲高}
D epoch 5のObliqueはNadirに比べMAE/RMSEが高く、local14は0.08657 m対0.03649 mであった。斜め視で地面平面の深度勾配、葉群の自己遮蔽、列間の見え方、遠近、撮影配置等が影響し得るが、ここでは仮説である。方向別Pearsonが絶対指標の順位と異なることも深度レンジと混合集計を検討する必要性を示す。
対象作物はCorn/Oat/Wheat/Soybeanで、稲を含まない。GTもカメラZ-depthで、植物頂部・地面の分割や地上座標鉛直高計算をしていない。したがってここでの結果は稲高精度を裏付けず、農作物画像への深度適応に関する予備的な結果に限る。
